%%%%%
%
\section{Moyenne-Echantillonnage d'un angle, avec adaptation � l'unit�}

\itemstyle{Op\'erateur~:} {\sf MoyEch\_Angl\_auto\_unit}
\index{MoyEch\_Angl\_auto\_unit}

%%%%%OBJET%%%%
\itemstyle{Objet~:}
Moyenne glissante et \'echantillonnage d'un param\`etre 
angulaire \`a une fr\'equence donn\'ee, diviseur de la fr\'equence du
param\`etre d'entr\'ee. Le param�tre d'entr�e peut �tre exprim�
en radians ou en degr�s.

%%%%%ARGUMENTS%%%%
\bigskip
%%%%%ENTREES%%%%
\begin{tabarg}
\multicolumn{3}{l}{\titretab {Entr\'ees}}\\
\hline
$X$ & Vecteur, \texttt{float} & Param\`etre d'entr\'ee ($rad$)\\
$R_f$ & Coeff., \texttt{double} & Rapport des fr\'equences ($ \frac{F_{sortie}}
{F_{entree}}$, dans l'intervalle $]0, 1]$)\\
$Decal$ & Coeff., \texttt{Int\_32} &  D\'ecalage (\textit{nombre de points},
 dans l'intervalle $[-1/R_f, 1/R_f]$)\\
$unite$ & Coeff., \texttt{String} & Unit� de $X$ (texte)\\
\end{tabarg}

\medskip
%%%%%SORTIES%%%%
\begin{tabarg}
\multicolumn{3}{l}{\titretab {Sorties}}\\
\hline
$X_m$ & Vecteur, \texttt{float} & R\'esultat (m�me unit� que $X$)\\
\end{tabarg}


%%%%%FORMULE%%%%
\itemstyle{Formule}

%%MoyEch_Angl_auto_unit
\[
Xm[i] = R_f . \sum_{j = i/R_f + Decal}^{(i+1)/R_f -1 + Decal} X[j]
\qquad i=0,Nbre\_Echant(X). R_f
\]
avec $|X[j+1] - X[j]| \leq (\pi \text{ ou } 180)$



%%%%COMMENTAIRE
\itemstyle{Commentaire}
Le nombre d'\'el\'ements du vecteur de sortie est le produit du
nombre d'\'echantillons du vecteur d'entr\'ee par $R_f$, arrondi
\`a l'entier sup\'erieur.

%%%%%AUTEUR%%%%
\itemstyle{Auteur~:}
CNRM/GMEI/TRAMM

%%%%%REFERENCE%%%%
\itemstyle{R\'ef\'erences~:}
